\documentclass[aspectratio=169]{beamer}

\usetheme{Madrid}
\usecolortheme{default}
\usefonttheme{professionalfonts}

\usepackage{amsmath,amssymb,bm}
\usepackage{braket}

\title{Information Entropy and von Neumann Entropy}
\subtitle{From Classical Probability to Quantum Density Matrices}
\author{}
\institute{}
\date{}

\begin{document}

% ==========================================================
\begin{frame}
  \titlepage
\end{frame}

% ==========================================================
\begin{frame}{Motivation: Quantifying uncertainty}
In both classical and quantum physics we need a quantitative measure of
\emph{uncertainty} or \emph{information content}.

\medskip
\begin{itemize}
\item Classical systems: uncertainty arises from probabilistic outcomes.
\item Quantum systems: uncertainty arises from superposition, measurement,
      and entanglement.
\end{itemize}

\medskip
Entropy provides a rigorous way to quantify:
\begin{itemize}
\item lack of information,
\item mixedness of a state,
\item correlations between subsystems.
\end{itemize}
\end{frame}

% ==========================================================
\begin{frame}{Classical information entropy (Shannon entropy)}
Consider a discrete random variable with outcomes $i$ occurring with
probabilities $\{p_i\}$, where
\[
p_i \ge 0, \qquad \sum_i p_i = 1.
\]

\medskip
The \textbf{Shannon entropy} is defined as
\[
H(\{p_i\}) = -\sum_i p_i \log_2 p_i.
\]

\medskip
Properties:
\begin{itemize}
\item $H \ge 0$.
\item $H=0$ if and only if one outcome has probability 1.
\item $H$ is maximal for a uniform distribution.
\end{itemize}

\medskip
For $N$ equally likely outcomes ($p_i=1/N$):
\[
H = \log_2 N.
\]
\end{frame}

% ==========================================================
\begin{frame}{Interpretation of Shannon entropy}
Shannon entropy quantifies:
\begin{itemize}
\item average information gained upon learning the outcome,
\item minimal number of bits needed to encode outcomes optimally,
\item uncertainty prior to measurement.
\end{itemize}

\medskip
Example:
\begin{itemize}
\item Fair coin: $p=(1/2,1/2)$ $\Rightarrow H=1$ bit.
\item Biased coin: $p=(0.9,0.1)$ $\Rightarrow H<1$ bit.
\end{itemize}

\medskip
\textbf{Key idea:} entropy depends only on probabilities, not on the
physical nature of the system.
\end{frame}

% ==========================================================
\begin{frame}{From classical probabilities to quantum states}
In quantum mechanics, measurement outcomes are probabilistic, but
the system is described by a state:

\medskip
\begin{itemize}
\item Pure state: $\ket{\psi}$.
\item Mixed state: classical uncertainty over quantum states.
\end{itemize}

\medskip
The appropriate object is the \textbf{density matrix}:
\[
\rho = \sum_i p_i \ket{\psi_i}\bra{\psi_i}.
\]

\medskip
Properties of $\rho$:
\begin{itemize}
\item $\rho^\dagger = \rho$ (Hermitian),
\item $\rho \ge 0$ (positive semidefinite),
\item $\Tr(\rho)=1$.
\end{itemize}
\end{frame}

% ==========================================================
\begin{frame}{Pure vs mixed states (density matrix viewpoint)}
\textbf{Pure state:}
\[
\rho = \ket{\psi}\bra{\psi}, \qquad \rho^2 = \rho, \qquad \Tr(\rho^2)=1.
\]

\medskip
\textbf{Mixed state:}
\[
\rho = \sum_i p_i \ket{\psi_i}\bra{\psi_i}, \qquad \rho^2 \neq \rho,
\qquad \Tr(\rho^2)<1.
\]

\medskip
Example (maximally mixed qubit):
\[
\rho = \frac{I_2}{2}
= \begin{pmatrix}
1/2 & 0\\
0 & 1/2
\end{pmatrix}.
\]

\medskip
This state represents complete ignorance about the qubit.
\end{frame}

% ==========================================================
\begin{frame}{Motivation for quantum entropy}
We seek a quantum analogue of Shannon entropy that:
\begin{itemize}
\item reduces to Shannon entropy for classical probability mixtures,
\item depends only on the density matrix $\rho$,
\item is invariant under unitary transformations,
\item quantifies mixedness and entanglement.
\end{itemize}

\medskip
This leads uniquely to the \textbf{von Neumann entropy}.
\end{frame}

% ==========================================================
\begin{frame}{Spectral decomposition of the density matrix}
Any density matrix can be diagonalized:
\[
\rho = \sum_i \lambda_i \ket{i}\bra{i},
\]
where:
\begin{itemize}
\item $\lambda_i \ge 0$ are eigenvalues,
\item $\sum_i \lambda_i = 1$,
\item $\{\ket{i}\}$ form an orthonormal basis.
\end{itemize}

\medskip
In this basis, $\rho$ behaves like a classical probability distribution
over orthogonal states.

\medskip
This observation motivates the definition of quantum entropy.
\end{frame}

% ==========================================================
\begin{frame}{Definition of von Neumann entropy}
The \textbf{von Neumann entropy} of a quantum state $\rho$ is
\[
S(\rho) = -\Tr\big(\rho \log_2 \rho\big).
\]

\medskip
Using the spectral decomposition:
\[
S(\rho) = -\sum_i \lambda_i \log_2 \lambda_i.
\]

\medskip
Thus, von Neumann entropy is precisely the Shannon entropy of the
eigenvalue distribution of $\rho$.
\end{frame}

% ==========================================================
\begin{frame}{Derivation: from density matrix to entropy}
Starting point:
\[
S(\rho) = -\Tr(\rho \log \rho).
\]

Diagonalize $\rho$:
\[
\rho = U \, \mathrm{diag}(\lambda_1,\dots,\lambda_n)\, U^\dagger.
\]

Then
\[
\log \rho = U \, \mathrm{diag}(\log \lambda_1,\dots,\log \lambda_n)\, U^\dagger.
\]

Using cyclicity of the trace:
\[
\Tr(\rho \log \rho)
= \Tr\big(\mathrm{diag}(\lambda_i)\,\mathrm{diag}(\log \lambda_i)\big)
= \sum_i \lambda_i \log \lambda_i.
\]

Therefore:
\[
S(\rho) = -\sum_i \lambda_i \log_2 \lambda_i.
\]
\end{frame}

% ==========================================================
\begin{frame}{Properties of von Neumann entropy}
\begin{itemize}
\item $S(\rho)\ge 0$.
\item $S(\rho)=0$ if and only if $\rho$ is pure.
\item $S(\rho)$ is maximal for $\rho=I/d$:
\[
S_{\max} = \log_2 d.
\]
\item Invariant under unitary evolution:
\[
S(U\rho U^\dagger) = S(\rho).
\]
\end{itemize}

\medskip
Entropy measures intrinsic mixedness, not basis-dependent uncertainty.
\end{frame}

% ==========================================================
\begin{frame}{Connection to Bell states and entanglement}
Consider the Bell state $\ket{\Phi^+}$:
\[
\ket{\Phi^+}=\frac{\ket{00}+\ket{11}}{\sqrt{2}}.
\]

\medskip
Total state is pure:
\[
\rho_{AB}=\ket{\Phi^+}\bra{\Phi^+}, \qquad S(\rho_{AB})=0.
\]

\medskip
Reduced state of one qubit:
\[
\rho_A = \Tr_B(\rho_{AB})=\frac{I_2}{2}.
\]

\medskip
Eigenvalues: $\lambda_1=\lambda_2=1/2$ $\Rightarrow$
\[
S(\rho_A)=1 \text{ bit}.
\]

\medskip
\textbf{Key insight:} entanglement produces local mixedness.
\end{frame}

% ==========================================================
\begin{frame}{Entropy as an entanglement indicator (pure states)}
For a bipartite pure state $\ket{\psi_{AB}}$:
\[
S(\rho_A)=S(\rho_B).
\]

\medskip
\begin{itemize}
\item $S(\rho_A)=0$ $\Rightarrow$ product (unentangled) state.
\item $S(\rho_A)>0$ $\Rightarrow$ entangled state.
\item Bell states: maximal entanglement, $S=1$ bit.
\end{itemize}

\medskip
Von Neumann entropy of subsystems provides a quantitative measure of
entanglement for pure bipartite states.
\end{frame}

% ==========================================================
\begin{frame}{Summary}
\begin{itemize}
\item Shannon entropy quantifies classical uncertainty.
\item Density matrices unify classical and quantum probabilities.
\item Von Neumann entropy:
\[
S(\rho)=-\Tr(\rho\log\rho),
\]
generalizes Shannon entropy to quantum mechanics.
\item For Bell states:
  \begin{itemize}
  \item global state is pure ($S=0$),
  \item subsystems are maximally mixed ($S=1$).
  \end{itemize}
\end{itemize}

\medskip
Entropy is a central concept linking information theory,
quantum mechanics, and entanglement.
\end{frame}

\end{document}


\documentclass[aspectratio=169]{beamer}

\usetheme{Madrid}
\usecolortheme{default}
\usefonttheme{professionalfonts}

\usepackage{amsmath,amssymb,bm}
\usepackage{braket}

\title{Quantum Entropy Beyond von Neumann}
\subtitle{Rényi Entropies, Measurement Entropy, and Mutual Information}
\author{}
\institute{}
\date{}

\begin{document}

% ==========================================================
\begin{frame}
  \titlepage
\end{frame}

% ==========================================================
\begin{frame}{Why go beyond von Neumann entropy?}
Von Neumann entropy captures intrinsic mixedness of a quantum state,
but additional entropy notions are useful to:
\begin{itemize}
\item characterize different aspects of uncertainty,
\item analyze measurements and classical outcomes,
\item quantify total vs purely quantum correlations,
\item study scaling and universality in many-body systems.
\end{itemize}

\medskip
This motivates extensions such as Rényi entropies,
measurement entropy, and mutual information.
\end{frame}

% ==========================================================
\section{Rényi Entropies}

\begin{frame}{Definition of Rényi entropy}
For a density matrix $\rho$, the \textbf{Rényi entropy} of order $\alpha$
is defined as
\[
S_\alpha(\rho)
= \frac{1}{1-\alpha} \log_2 \Tr(\rho^\alpha),
\qquad \alpha>0,\ \alpha\neq 1.
\]

\medskip
Special cases:
\begin{itemize}
\item $\alpha \to 1$: von Neumann entropy
\[
\lim_{\alpha\to 1} S_\alpha(\rho)=S(\rho).
\]
\item $\alpha=2$: second Rényi entropy
\[
S_2(\rho)=-\log_2 \Tr(\rho^2).
\]
\end{itemize}
\end{frame}

% ==========================================================
\begin{frame}{Interpretation of Rényi entropies}
Let $\{\lambda_i\}$ be eigenvalues of $\rho$.

\medskip
\[
S_\alpha(\rho)
= \frac{1}{1-\alpha} \log_2 \sum_i \lambda_i^\alpha.
\]

\medskip
Interpretation:
\begin{itemize}
\item Larger $\alpha$: emphasizes largest eigenvalues.
\item Smaller $\alpha$: sensitive to full spectrum.
\item $S_2$ directly related to purity $\Tr(\rho^2)$.
\end{itemize}

\medskip
Rényi entropies are widely used in:
\begin{itemize}
\item numerical many-body physics,
\item entanglement scaling and area laws,
\item experimental protocols (e.g. swap tests).
\end{itemize}
\end{frame}

% ==========================================================
\begin{frame}{Rényi entropy for Bell states}
For a Bell state $\ket{\Phi^+}$:
\[
\rho_A = \frac{I_2}{2}.
\]

\medskip
Eigenvalues: $\lambda_1=\lambda_2=1/2$.

\medskip
Thus, for any $\alpha$:
\[
S_\alpha(\rho_A)
= \frac{1}{1-\alpha}\log_2\left(2\cdot(1/2)^\alpha\right)
= 1.
\]

\medskip
\textbf{Conclusion:}
Bell states are maximally entangled according to all Rényi entropies.
\end{frame}

% ==========================================================
\section{Measurement Entropy vs State Entropy}

\begin{frame}{Measurement entropy}
Given a state $\rho$ and a projective measurement
$\{ \Pi_i \}$, the outcome probabilities are
\[
p_i = \Tr(\rho\,\Pi_i).
\]

\medskip
The \textbf{measurement entropy} is the Shannon entropy
\[
H_{\text{meas}} = -\sum_i p_i \log_2 p_i.
\]

\medskip
This entropy depends on:
\begin{itemize}
\item the quantum state $\rho$,
\item the choice of measurement basis.
\end{itemize}
\end{frame}

% ==========================================================
\begin{frame}{State entropy vs measurement entropy}
\textbf{Von Neumann entropy}:
\[
S(\rho) = -\Tr(\rho \log \rho)
\]
\begin{itemize}
\item intrinsic property of the state,
\item basis independent.
\end{itemize}

\medskip
\textbf{Measurement entropy}:
\[
H_{\text{meas}}(\{\Pi_i\}) = -\sum_i p_i\log p_i
\]
\begin{itemize}
\item depends on measurement choice,
\item reflects classical uncertainty of outcomes.
\end{itemize}

\medskip
General inequality:
\[
H_{\text{meas}} \ge S(\rho),
\]
with equality if measurement is in the eigenbasis of $\rho$.
\end{frame}

% ==========================================================
\begin{frame}{Example: Bell state measurement}
For $\ket{\Phi^+}$:
\[
\rho_A=\frac{I_2}{2}, \qquad S(\rho_A)=1.
\]

\medskip
Measuring $\sigma_z$:
\[
p(0)=p(1)=\frac{1}{2}
\quad \Rightarrow \quad H_{\text{meas}}=1.
\]

\medskip
Measuring $\sigma_x$:
\[
p(+)=p(-)=\frac{1}{2}
\quad \Rightarrow \quad H_{\text{meas}}=1.
\]

\medskip
Here $H_{\text{meas}}=S(\rho)$ for all bases
because $\rho_A$ is maximally mixed.
\end{frame}

% ==========================================================
\section{Mutual Information and Entanglement}

\begin{frame}{Quantum mutual information}
For a bipartite state $\rho_{AB}$, the
\textbf{quantum mutual information} is defined as
\[
I(A:B)
= S(\rho_A) + S(\rho_B) - S(\rho_{AB}).
\]

\medskip
It quantifies:
\begin{itemize}
\item total correlations between $A$ and $B$,
\item both classical and quantum correlations.
\end{itemize}

\medskip
Always non-negative:
\[
I(A:B)\ge 0.
\]
\end{frame}

% ==========================================================
\begin{frame}{Mutual information for Bell states}
For a Bell state $\ket{\Phi^+}$:
\[
S(\rho_{AB})=0,
\quad
S(\rho_A)=S(\rho_B)=1.
\]

\medskip
Thus:
\[
I(A:B) = 1+1-0 = 2.
\]

\medskip
Interpretation:
\begin{itemize}
\item maximal total correlation,
\item all correlations are quantum in origin.
\end{itemize}
\end{frame}

% ==========================================================
\begin{frame}{Entanglement entropy}
For a pure bipartite state $\ket{\psi_{AB}}$,
the \textbf{entanglement entropy} is defined as
\[
S_E = S(\rho_A)=S(\rho_B).
\]

\medskip
Properties:
\begin{itemize}
\item $S_E=0$ for product states,
\item $S_E>0$ for entangled states,
\item maximal for Bell states.
\end{itemize}

\medskip
For mixed states, entanglement entropy alone is insufficient;
more refined measures are required.
\end{frame}

% ==========================================================
\begin{frame}{Many-body perspective}
In many-body quantum systems:
\begin{itemize}
\item entanglement entropy scales with subsystem size,
\item ground states often obey area laws,
\item critical systems show logarithmic corrections.
\end{itemize}

\medskip
Rényi entropies are especially useful in:
\begin{itemize}
\item tensor-network simulations,
\item quantum Monte Carlo,
\item experimental cold-atom platforms.
\end{itemize}
\end{frame}

% ==========================================================
\begin{frame}{Summary}
\begin{itemize}
\item Rényi entropies generalize von Neumann entropy and probe
      different parts of the spectrum.
\item Measurement entropy quantifies classical uncertainty
      induced by a measurement.
\item Quantum mutual information measures total correlations.
\item Bell states:
  \begin{itemize}
  \item zero global entropy,
  \item maximal local entropy,
  \item maximal mutual information.
  \end{itemize}
\end{itemize}

\medskip
Together, these concepts form the core entropy toolkit
of quantum information theory.
\end{frame}

\end{document}
